% ---- pg-032 （书页 20） ----

\begin{mvfigure}
\mvinc{pg-032}{0.147,0.125,0.827,0.403}
\phantomsection\label{fig:1.15}\mvcap{图 1.15\quad 单个母灰熊身长的生长曲线列}
\end{mvfigure}

\mvsub{星形}

假设每个数据单元由 $p\geqslant 2$ 个变量的非负观测值组成.在二维空间中,我们可构造具有固定(参照)半径的圆,从圆心引出 $p$ 条等距的射线,这些射线的长度代表变量的数值,以直线连接射线的端点即形成一个星形.每个星形代表一个多元观测值,这些星形可根据它们(主观的)的相似性分组.

在构造星形时,将观测值标准化常常会有所帮助.在这种情况下,一些观测值会是负数.观测值可以重新表示,使得圆心代表整个数据集中的最小标准化观测值.

\mvbold{例 1.11}\quad(星形的实用数据)

\hyperref[fig:1.16]{图 1.16} 给出 5 个星形图.它们代表第 12 章\hyperref[tab:12.4]{表 12.4} 所列出的 22 个公用事业公司中的前 5 个公司的观测值.共有 8 个变量,因此,这些星形是变形的八边形.

所有变量的观测值都是经过标准化的.在这前 5 个公司中,对任何变量最小的标准化观测值是 $-1.6$.将这个值作为零点,沿着从圆心引出的 8 条等角射线以相同的比例画出 8 个变量,这些变量以 12 点钟的位置为起点,按顺时针方向排序.

初看起来,这些公司中没有一个与其他任何一个是相似的.由于星形的构造方法,不管怎样,每个变量在感观上权数是相等的.如果我们仔细观察变量 6(以每年使用千瓦时为单位的销售额总量)和变量 8(以美分/千瓦时为单位的燃料费用),则发现,波士顿爱迪生和联合爱迪生是相似的(小变量 6,大变量 8),以及亚利桑那公共服务公司、中央路易斯安那电力公司和联邦爱迪生公司是相似的(中等变量 6 和中等变量 8).

\mvsub{切尔诺夫脸}

人们的反应表现在脸上.切尔诺夫(见文献[4])假定用二维平面的脸来表示 $p$ 维观测结果,脸的特征(如脸的形状,嘴的弯曲率,鼻子的长度,眼睛的大小,瞳孔的位置等等)是由 $p$ 个变量的测量值所决定的.

按照最初的设计,切尔诺夫脸可处理多达 18 个变量.脸部容貌对应的变量的分配是由实%
